% CVPR 2026 Paper Template; see https://github.com/cvpr-org/author-kit

\documentclass[10pt,twocolumn,letterpaper]{article}

%%%%%%%%% PAPER TYPE  - PLEASE UPDATE FOR FINAL VERSION
\usepackage{cvpr}              % To produce the CAMERA-READY version
% \usepackage[review]{cvpr}      % To produce the REVIEW version
% \usepackage[pagenumbers]{cvpr} % To force page numbers, e.g. for an arXiv version

\usepackage{cuted}
\usepackage{stfloats}
\usepackage{placeins}
\usepackage{multirow} 


% Import additional packages in the preamble file, before hyperref
\input{preamble}
\input{macros}
\input{macro} % macros Beibei

% It is strongly recommended to use hyperref, especially for the review version.
% hyperref with option pagebackref eases the reviewers' job.
% Please disable hyperref *only* if you encounter grave issues, 
% e.g. with the file validation for the camera-ready version.
%
% If you comment hyperref and then uncomment it, you should delete *.aux before re-running LaTeX.
% (Or just hit 'q' on the first LaTeX run, let it finish, and you should be clear).
\definecolor{cvprblue}{rgb}{0.21,0.49,0.74}
\usepackage[pagebackref,breaklinks,colorlinks,allcolors=cvprblue]{hyperref}

%%%%%%%%% PAPER ID  - PLEASE UPDATE
\def\paperID{9517} % *** Enter the Paper ID here
\def\confName{CVPR}
\def\confYear{2026}

%%%%%%%%% TITLE - PLEASE UPDATE
% \title{\LaTeX\ Author Guidelines for \confName~Proceedings}
\title{FabricGen: Microstructure-Aware Woven Fabric Generation}


%%%%%%%%% AUTHORS - PLEASE UPDATE
% For a paper whose authors are all at the same institution,
% omit the following lines up until the closing ``}''.
% Additional authors and addresses can be added with ``\and'',
% just like the second author.
% To save space, use either the email address or home page, not both
% \author{Yingjie Tang\\
% Nankai University\\
% {\tt\small lxtyin.ac@gmail.com}
% \and
% Di Luo\\
% Nankai University\\
% {\tt\small diluo@mail.nankai.edu.cn}
% \and
% Zixiong Wang\\
% Nankai University\\
% {\tt\small zixiong_wang@outlook.com}
% \and
% Xiaoli Ling\\
% School of Intelligence Science and Technology, Nanjing University\\
% {\tt\small xiaoliling@smail.nju.edu.cn}
% \and
% Jian Yang\\
% Nankai University\\
% {\tt\small csjyang@nankai.edu.cn}
% \and
% Beibei Wang\\
% School of Intelligence Science and Technology, Nanjing University\\
% {\tt\small beibei.wang@nju.edu.cn}
% }

\author{
    Yingjie Tang$^{1\dagger}$ \quad 
    Di Luo$^{1\dagger}$ \quad 
    Zixiong Wang$^{1}$ \quad 
    Xiaoli Ling$^{2}$ \quad 
    Jian Yang$^{1,2}$ \quad 
    Beibei Wang$^{2*}$ \\
    {\normalsize $^{1}$Nankai University} \quad {\normalsize $^{2}$Nanjing University}
}

\usepackage{mathtools}
\usepackage{xspace}


\begin{document}
\maketitle

\begin{strip}
    \centering
    \vspace{-40px}
    \includegraphics[width=1.0\textwidth]{fig/teaser3.pdf}
    \captionof{figure}{
We present FabricGen, an end-to-end \revise{material generation} framework that generates woven fabric materials with fine details, enabling non-expert users to create diverse and highly detailed fabric materials for photorealistic rendering. 
    }
\label{fig:teaser}
\end{strip}

\newcommand\blfootnote[1]{%
  \begingroup
  \renewcommand\thefootnote{}\footnote{#1}%
  \addtocounter{footnote}{-1}%
  \endgroup
}
\blfootnote{* Corresponding author.}
\blfootnote{$\dagger$ These authors contribute equally.}
\blfootnote{$^1$ College of Computer Science, Nankai University, Tianjin, China}
\blfootnote{$^2$ School of Intelligence Science and Technology, Nanjing University, Suzhou, China}



% \input{sec/0_abstract}    
% \input{sec/1_intro}
% \input{sec/2_formatting}
% \input{sec/3_finalcopy}
\input{0_sec_abstract}
\input{1_sec_intro}
\input{2_sec_related}
\input{3_sec_background}
\input{4_sec_framework}
\input{5_sec_procedural}
\input{6_sec_results}
\input{7_sec_conclusion}

\section*{Acknowledgment}
We thank the reviewers for the valuable comments. This work has been partially supported by the National Natural Science Foundation of China under grant No. 62572230.

%-----------------------------------%
% Figure only page
%-----------------------------------%

% \begin{figure}[tb]
% \centering
% \includegraphics[width = 0.75\linewidth]{./fig/placeholder.pdf}
% \caption*{place holder}
% \end{figure}

% \begin{figure}[h]
% \centering
% \includegraphics[width = 0.95\linewidth]{./fig/ablation_flyaway.pdf}
% \caption{ The impact of flyaway fibers. Fibers escaping from the surface cause irregular highlights in real fabrics. While such effects are ignored in previous methods~\cite{Jin:2022:inverse}, our approach accurately captures this behavior. }
% \label{fig:ablation_flyaway}
% \end{figure}

\clearpage
% \FloatBarrier  % 确保在 reference 之前强制输出
\onecolumn
\twocolumn


{
    \small
    \bibliographystyle{ieeenat_fullname}
    \bibliography{main}
}

% WARNING: do not forget to delete the supplementary pages from your submission 
% \input{sec/X_suppl}

\end{document}
